\subsection{汴州的朝代与江淮的水网}
907年以后，梁、唐、晋、汉、周在中原相继建立；同一时期的江淮、两浙、湖湘、岭南、闽岭、巴蜀和河东并没有暂停自己的政治时间。汴州能够成为中心，靠的是黄河、运河、仓储和军镇的交会；它也因此特别容易受北方军队、河道与财政中断的牵制。把五代写成五个国号的顺序，会看不见城市、渡口和军粮如何使一个政权得到或失去行动余地。

923年后唐进入汴州，见\eventref{LT-923-092}{后唐灭后梁}。这既是军事胜负，也是接收一座都城及其官署、仓场、旧部与运输联系的过程。正史能较好地保留称帝、入城和任命的顺序，较难告诉我们每一处市场、乡村与军营怎样调整。中原政权更替频繁，并不表示每次更替都以相同强度到达州县；地方官、军人、商人和宗族常在新的名义下继续处理粮、税和安全。

吴与后来南唐所处的江淮水网则带来另一种条件。杨氏在907年仍用唐号，见\eventref{LT-907-084}{杨渥续用唐号}，说明旧唐的合法性语言可以被地方政权利用，而不等于它受开封或长安的直接指挥。运河支流、盐业、稻作、城市与江面运输使江淮能够筹集资源，也使控制水路、仓场和地方中介成为持续的政治工作。钱币、墓志和地方记述能补出少数节点，却不能把南方的经济活力换算成每个家庭的生活。

\subsection{杭州、成都、广州与太原各有其边界}
吴越面对的是钱塘江、杭州城、浙东水路和海上交往。它与中原政权保持册封、贡使和外交关系时，所使用的称号不应被读成行政从属的精确地图；使者、商船、寺院与地方官各自的路线使这种关系具有实际内容。吴越的水利、城防和佛教赞助也不能只当作“偏安”的装饰，它们是区域统治在水网和人口密集地带寻找稳定的方式。

前后蜀在成都平原与西南山路之间经营。巴蜀的丰饶并非一个自然自动提供的金库：山口、河流、地方军队、盐与粮食的转运都决定资源能否被都城调动。中原战乱使入蜀与出蜀的路径具有避难与军事双重意义，但不能把每一次人口移动都想象为同样方向、同样规模的“南迁”。墓志和地方遗存只能显示可保存的少数轨迹，正因此需要把不可见的大多数留在判断的边界之外。

南汉的广州与岭南海路又连接着不同的市场和外来信息。港口遗存、陶瓷、钱币与后出的记载提示贸易的存在，却不能单凭一船货物或一处窖藏断言国家直接控制了全部海上交换。海路给地方政权带来税源、物资和外交机会，也使风季、航道、商人网络与沿岸社会成为不可忽略的约束。把岭南只写作中原王朝的最南端，会失去它作为海上节点的行动能力。

闽、楚、荆南与北汉的处境同样不能折叠。闽岭的山道与港口、湖南的湖沼和粮运、荆南卡在江汉交通的位置、太原与契丹及北方通道的关系，分别塑造了能动员什么、需向谁谈判、最怕哪条路被切断。十国国史散佚严重，后来的汇编与中原史书带有强烈选择；这不是可以用想象补满的缺口，而是要求叙事承认不同区域的可见度并不相等。

\subsection{契丹援军、燕云与后周的选择}
936年石敬瑭在契丹援助下建立后晋，见\eventref{LT-936-099}{后晋建国}。把这一事件只写成“割让燕云”的道德故事，会遮去太原起兵、北方骑兵通道、城镇与新政权合法性之间的交换。现存叙事对盟约细节和边界多有后设整理，故不能假装能画出每座城在每一年归属何方；但燕云一带的关口与道路为何会改变中原防御和契丹的战略选择，则是可以追问的空间问题。

后周的整军、财政与对南方的行动，显示中原并非没有恢复调度能力。它的改革能否落实，仍取决于军队补给、州县征收、铸钱、河运与既有军镇的合作。960年赵匡胤取代后周，给下一卷提供了新的中枢，却没有把五代十国的区域积累一笔勾销。宋的统一是在这些道路、市场、军政集团与多政权交涉的遗产上继续发生的；读者离开本卷时，应把“统一”再一次理解为接收、谈判和重新分配成本的过程。
